
\documentclass{article}
\usepackage{colortbl}
\usepackage{makecell}
\usepackage{multirow}
\usepackage{supertabular}

\begin{document}

\newcounter{utterance}

\centering \large Interaction Transcript for game `cladder', experiment `full\_v1.5\_default', episode 7000 with Qwen3.6{-}27B.
\vspace{24pt}

{ \footnotesize  \setcounter{utterance}{1}
\setlength{\tabcolsep}{0pt}
\begin{supertabular}{c@{$\;$}|p{.15\linewidth}@{}p{.15\linewidth}p{.15\linewidth}p{.15\linewidth}p{.15\linewidth}p{.15\linewidth}}
    \# & \multicolumn{2}{c}{Player} && \multicolumn{2}{c}{Game Master} \\
    \hline

    \theutterance \stepcounter{utterance}  
    & & & \multicolumn{4}{p{0.6\linewidth}}{
        \cellcolor[rgb]{0.9,0.9,0.9}{
            \makecell[{{p{\linewidth}}}]{
                \texttt{\tiny{[P1$\langle$GM]}}
                \texttt{You are an expert in causal inference. The following question is not a typical commonsense query, but rather a meticulously designed question created by a professor specializing in causal inference, intended to assess the students' mastery of the course content.} \\
\texttt{Imagine a self{-}contained, hypothetical world with only the following conditions, and without any unmentioned factors or causal relationships: Jyka has a direct effect on gyzp and hwax. Hwax has a direct effect on lirg. Gyzp has a direct effect on lirg. For those who are not jyka, the probability of lirg is 63\%. For those who are jyka, the probability of lirg is 49\%. For those who are jyka, would it be more likely to see lirg if the individual was not jyka?} \\
\texttt{Start your answer with 'Yes' or 'No', followed by additional reasoning or evidence to support your explanation.} \\
            }
        }
    }
     \\ \\

    \theutterance \stepcounter{utterance}  
    & \multicolumn{4}{p{0.6\linewidth}}{
        \cellcolor[rgb]{0.561,0.737,0.561}{
            \makecell[{{p{\linewidth}}}]{
                \texttt{\tiny{[P1$\rangle$GM]}}
                \texttt{Yes} \\
\\ 
\texttt{**Reasoning:**} \\
\\ 
\texttt{To answer this question, we must distinguish between the **observed probability** (association) and the **counterfactual probability** (causal effect). The question asks: "For those who are jyka, would it be more likely to see lirg if the individual was not jyka?" This is a counterfactual query asking for the probability of the outcome \$Lirg\$ under the intervention or hypothetical scenario where \$Jyka\$ is set to 0 (not jyka), specifically for the subpopulation that actually has \$Jyka=1\$.} \\
\\ 
\texttt{Let's break down the causal structure and probabilities provided:} \\
\\ 
\texttt{1.  **Causal Structure:**} \\
\texttt{    *   \$Jyka \textbackslash{}rightarrow Gyzp\$} \\
\texttt{    *   \$Jyka \textbackslash{}rightarrow Hwax\$} \\
\texttt{    *   \$Hwax \textbackslash{}rightarrow Lirg\$} \\
\texttt{    *   \$Gyzp \textbackslash{}rightarrow Lirg\$} \\
\texttt{    *   There are no other factors. This is a Directed Acyclic Graph (DAG).} \\
\\ 
\texttt{2.  **Observed Probabilities:**} \\
\texttt{    *   \$P(Lirg | Jyka = 0) = 0.63\$} \\
\texttt{    *   \$P(Lirg | Jyka = 1) = 0.49\$} \\
\\ 
\texttt{3.  **The Counterfactual Question:**} \\
\texttt{    The question asks for \$P(Lirg\_\{Jyka=0} \\
            }
        }
    }
    & & \\ \\

    \theutterance \stepcounter{utterance}  
    & & & \multicolumn{2}{p{0.3\linewidth}}{
        \cellcolor[rgb]{0.9,0.9,0.9}{
            \makecell[{{p{\linewidth}}}]{
                \texttt{\tiny{[GM$|$GM]}}
                \texttt{Yes} \\
            }
        }
    }
    & & \\ \\

    \theutterance \stepcounter{utterance}  
    & & & \multicolumn{2}{p{0.3\linewidth}}{
        \cellcolor[rgb]{0.9,0.9,0.9}{
            \makecell[{{p{\linewidth}}}]{
                \texttt{\tiny{[GM$|$GM]}}
                \texttt{yes} \\
            }
        }
    }
    & & \\ \\

    \theutterance \stepcounter{utterance}  
    & & & \multicolumn{2}{p{0.3\linewidth}}{
        \cellcolor[rgb]{0.9,0.9,0.9}{
            \makecell[{{p{\linewidth}}}]{
                \texttt{\tiny{[GM$|$GM]}}
                \texttt{game\_result = WIN} \\
            }
        }
    }
    & & \\ \\

\end{supertabular}
}

\end{document}
